% ============================================================
%  Preàmbul compartit — Aprendes Universitat
%  Compartit pels llibres d'Aprendes Universitat.
%  Les dades pròpies de cada llibre (títol, autor i metadades PDF)
%  van al metadata.tex local de cada projecte, carregat ABANS
%  d'aquest fitxer. La macro \hypermetadata es crida DESPRÉS,
%  quan hyperref ja ha estat carregat.
% ============================================================

\usepackage{eurosym}
\usepackage{iftex}
\ifpdftex
  \usepackage[T1]{fontenc}
  \usepackage[utf8]{inputenc}
  \usepackage{lmodern}
\else
  % XeLaTeX / LuaLaTeX: Unicode natiu, res d'inputenc
  \usepackage{fontspec}
  % Carrega Latin Modern per nom de fitxer, que XeTeX/LuaTeX pot
  % resoldre via kpathsea fins i tot si fontconfig no l'ha registrat.
  \setmainfont{lmroman10-regular.otf}[
    BoldFont=lmroman10-bold.otf,
    ItalicFont=lmroman10-italic.otf,
    BoldItalicFont=lmroman10-bolditalic.otf
  ]
  \setsansfont{lmsans10-regular.otf}[
    BoldFont=lmsans10-bold.otf,
    ItalicFont=lmsans10-oblique.otf,
    BoldItalicFont=lmsans10-boldoblique.otf
  ]
  \setmonofont{lmmono10-regular.otf}[ItalicFont=lmmono10-italic.otf]
\fi
% microtype: protrusion i expansió completes només amb pdfLaTeX.
% Amb XeLaTeX/LuaLaTeX només hi ha protrusion i generaria avisos de
% «slot number»; per això només l'activem sota pdfTeX.
\ifpdftex
  \usepackage{microtype}
\fi
\usepackage[catalan, provide=*]{babel}
\usepackage[autostyle=true,babel=true]{csquotes}
% csquotes no porta estil català: el declarem perquè \enquote{...}
% doni guillemets «...» (exteriors) i cometes altes “...” (interiors),
% independentment de la versió de csquotes instal·lada.
\DeclareQuoteStyle{catalan}
  {\guillemotleft}{\guillemotright}
  {\quotedblbase}{\textquotedblleft}
\DeclareQuoteAlias{catalan}{catala}
\usepackage{amsmath,mathtools}% amssymb, amsthm: via aprendes.sty
\usepackage{geometry}
\usepackage{import}
\usepackage{graphicx}
\usepackage{tabularx}
\usepackage{tikz}
\usetikzlibrary{positioning,arrows.meta,matrix}
\usepackage[all]{xy}% diagrames de grafs (llaços i arestes múltiples) a l'apèndix de grafs
\usepackage{float}
\usepackage[section]{placeins}% FloatBarrier automàtic: cap figura no travessa una secció
\usepackage{pgfplots}
\pgfplotsset{compat=1.18}
\usepackage{imakeidx}
\usepackage{booktabs}

% La bibliografia es configura localment en cada projecte quan hi ha referències.

\setcounter{MaxMatrixCols}{10}

\graphicspath{{img/}}
\geometry{hmargin={3cm,3cm},vmargin={2.5cm,2.5cm},offset=0pt,footskip=1cm,headsep=1cm}


\setlist{itemsep=0.25em, topsep=0.25em}
\numberwithin{equation}{section}
\setcounter{secnumdepth}{3}
\setcounter{tocdepth}{2}

% Configuració dels paràgrafs
\usepackage{indentfirst}
\setlength{\parindent}{1.5em}
\setlength{\parskip}{0.3em}

% ============================================================
%  hyperref: ÚLTIM paquet, com recomana el seu manual.
%  Els colors dels enllaços es fixen aquí. Les metadades
%  pròpies del llibre ja s'han definit abans en metadata.tex;
%  la macro \hypermetadata s'executa després de carregar aquest fitxer.
% ============================================================
\usepackage{hyperref}
% Color dels enllaços: el blau pissarra de la paleta d'aprendes,
% una mica enfosquit. Prou distint del vermell corporatiu dels
% títols, exemples i exercicis, i coherent amb la resta de colors.
\colorlet{EnllacBlau}{ResultatBlau!90!black}
\hypersetup{
  colorlinks=true,
  linktoc=page,
  linkcolor=EnllacBlau,
  urlcolor=EnllacBlau,
  citecolor=EnllacBlau,
  filecolor=EnllacBlau,
  hypertexnames=false,
}
\usepackage{bookmark}% millora els marcadors del PDF; ha d'anar després d'hyperref
